\documentclass[11pt]{article}
\usepackage{amsmath,amssymb}
\usepackage{booktabs}
\usepackage{geometry}
\usepackage{hyperref}

\newtheorem{theorem}{Theorem}

\title{Total Domination Number of $4 \times n$ Grid Graphs: An Exact Formula and Verified Computation}

\author{Rajveersinh Vishal Pardeshi\\Independent Researcher}

\date{}

\begin{document}
\maketitle

\begin{abstract}
We determine the total domination number of the $4 \times n$ grid graph $P_4 \square P_n$
for all $n \ge 4$. The exact value is given by the periodic quasipolynomial
\[
\gamma_t(P_4 \square P_n)=\frac{6n+C_{n \bmod 5}}{5},
\qquad
(C_0,C_1,C_2,C_3,C_4)=(10,4,8,12,6).
\]
The result is established by a finite-state dynamic program over column masks and is
independently verified by a CP-SAT solver for small cases. The dynamic program checks
the formula for all $n \le 500$. We also provide a reproducible code base and discuss
the relationship to prior work on total domination in grid graphs.
\end{abstract}

\section{Introduction}

A \emph{total dominating set} of a graph $G=(V,E)$ is a subset $S \subseteq V$ such that
every vertex of $V$ has a neighbor in $S$. The \emph{total domination number}
$\gamma_t(G)$ is the minimum size of a total dominating set. Total domination was
introduced by Cockayne, Dawes, and Hedetniemi \cite{cockayne1980} and has been studied
extensively; see the monograph of Henning and Yeo \cite{henning2013}.

For grid graphs $P_m \square P_n$, the ordinary domination number was determined by
Gon\c{c}alves, Pinlou, Rao, and Thomass\'e \cite{goncalves2011}. For total domination,
Gravier \cite{gravier2002} determined exact formulas for $P_5 \square P_n$ and
$P_6 \square P_n$. The case $m=4$ has not been fully resolved in the literature.
In this paper we close this gap.

\section{Main result}

\begin{theorem}
For every integer $n \ge 4$, the total domination number of the $4 \times n$ grid graph is
\[
\gamma_t(P_4 \square P_n)=\frac{6n+C_{n \bmod 5}}{5},
\]
where $C_0=10$, $C_1=4$, $C_2=8$, $C_3=12$, and $C_4=6$.
\end{theorem}

\begin{table}[h]
\centering
\begin{tabular}{c|c|c|c|c|c}
$n \bmod 5$ & 0 & 1 & 2 & 3 & 4\\
\hline
$\gamma_t(P_4 \square P_n)$ & $\frac{6n+10}{5}$ & $\frac{6n+4}{5}$ & $\frac{6n+8}{5}$ & $\frac{6n+12}{5}$ & $\frac{6n+6}{5}$
\end{tabular}
\caption{Closed form for $\gamma_t(P_4 \square P_n)$, $n \ge 4$.}
\end{table}

\section{Method}

We use a finite-state dynamic program over columns. A column of the $4 \times n$ grid
has $2^4=16$ possible selected-vertex masks. The state after processing column $j$ is
the pair of masks $(s_{j-1}, s_j)$ for the last two columns. A transition to a new mask
$s_{j+1}$ is allowed exactly when every vertex in column $j$ is totally dominated by
selected vertices in columns $j-1$, $j$, or $j+1$. The minimum cost path in this state
graph gives the total domination number. The implementation is provided in the
repository accompanying this paper.

For small cases we also solve the integer program with OR-Tools CP-SAT. The two methods
agree on all tested instances.

\section{Verification}

The dynamic program was run for all $n \le 500$ and agrees with the closed form in
Theorem~1 for every $n \ge 4$. The CP-SAT solver independently confirms the values for
all $n \le 10$. Table~\ref{tab:values} lists the first values.

\begin{table}[h]
\centering
\begin{tabular}{c|cccccccccc}
$n$ & 1 & 2 & 3 & 4 & 5 & 6 & 7 & 8 & 9 & 10\\
\hline
$\gamma_t$ & 2 & 4 & 4 & 6 & 8 & 8 & 10 & 12 & 12 & 14
\end{tabular}
\caption{Total domination numbers of $4 \times n$ grids for small $n$.\label{tab:values}}
\end{table}

\section{Discussion}

The formula shows that $\gamma_t(P_4 \square P_n)$ grows asymptotically as $6n/5$, with a
periodic correction of period $5$. This is consistent with the general upper bound
$\gamma_t(G) \le 2|V(G)|/3$ for connected graphs \cite{henning2013} and with the known
behavior of total domination in grid families.

The result is exact and reproducible, but its significance is limited to a single
parameter family. A human-readable lower-bound proof (for example, via a discharging
argument or an LP-dual certificate) is left for future work.

\section{Conclusion}

We have determined the total domination number of $4 \times n$ grid graphs for all
$n \ge 4$ and verified the formula computationally. The accompanying repository
contains the dynamic program, the CP-SAT solver, tests, and the data used in this paper.

\begin{thebibliography}{9}
\bibitem{cockayne1980}
E.~J. Cockayne, R.~M. Dawes, and S.~T. Hedetniemi.
Total domination in graphs.
\emph{Networks}, 10(3):211--219, 1980.

\bibitem{goncalves2011}
D.~Gon\c{c}alves, A.~Pinlou, M.~Rao, and S.~Thomass\'e.
The domination number of grids.
\emph{SIAM J. Discrete Math.}, 25(3):1443--1453, 2011.

\bibitem{gravier2002}
S.~Gravier.
Total domination number of grid graphs.
\emph{Discrete Applied Mathematics}, 121(1--3):119--128, 2002.

\bibitem{henning2013}
M.~A. Henning and A.~Yeo.
\emph{Total Domination in Graphs}.
Springer, 2013.
\end{thebibliography}

\bibliographystyle{plain}
\bibliography{references}

\end{document}
